\section{Genuine traces, exact lengths, and testing}\label{sec:transfer}

Proposition~\ref{prop:construction} now supplies the final pair and its auxiliary laws.  We first retain a one-trace bound before using product observations.  This proves the minimum-total-variation theorem independently of the sample-complexity deduction.

\begin{lemma}[Restoration and padding]\label{lem:restoration}
For the parameters in Section~\ref{sec:scales}, there are distinct words $x,y\in\{0,1\}^n$ such that, for every positive integer $m$,
\begin{equation}\label{eq:restored-product}
 \TV\bigl(\Dcal_q(x)^{\otimes m},\Dcal_q(y)^{\otimes m}\bigr)
 \le \sqrt{m n^{-d_D}}+2mn(D+1)\Delta.
\end{equation}
In particular,
\begin{equation}\label{eq:restored-single}
 \TV\bigl(\Dcal_q(x),\Dcal_q(y)\bigr)
 \le n^{-d_D/2}+2n(D+1)\Delta.
\end{equation}
\end{lemma}

\begin{proof}
Append the same arbitrary binary tail to $x_D,y_D$ to reach length $n$.  They remain distinct.  Applying the final common oracle kernel, then generating an independent trace of the known tail and concatenating it, approximates each padded $\delta$-trace law within $E_D\le n(D+1)\Delta$ in total variation.  Common padding by this kernel is the elementary monotonicity principle also used by \citet[proof of Lemma 10]{MPV2014}.

The final moment estimate, monotonicity of norms on a probability space, and $r_D\ge2$ give
\[
 H^2(\mu_D,\nu_D)
 \le\|d\nu_D/d\mu_D-1\|_{L^2(\mu_D)}^2
 \le n^{-d_D}.
\]
Tensorization and $\TV\le H$ bound the oracle product distance by $\sqrt{m n^{-d_D}}$.  Replacing each of the $m$ projected factors by its genuine law costs at most $mE_D$ for each hypothesis, by tensor telescoping.  This proves \eqref{eq:restored-product} for $\delta$.  Applying the additional deletion kernel independently to each sample proves the same bound for the requested $q$.  Taking $m=1$ gives \eqref{eq:restored-single}.
\end{proof}

\begin{proof}[Proof of Theorem~\ref{thm:tv}]
Fix any $q\in(0,1)$.  Choose $B=2$, $\zeta=1/2$, and auxiliary coefficients $c=1/(80\log2)<c_+=1/(40\log2)$ once, independently of $A$.  The setup inequality holds because its left side is $1/(20\log2)$.  The construction remains valid with these fixed parameters; its proof requires only their positivity and the stated eventual inequalities.  In this case $\delta>0$ is fixed, $D\to\infty$, and $d_D=2.20+0.20D\to\infty$.  For every fixed $A>0$, multiplying \eqref{eq:restored-single} by $n^A$ gives
\[
 n^{A-d_D/2}+2(D+1)n^{A+1-100(d_D+B+1)}\longrightarrow0.
\]
The construction is available for every sufficiently large integer $n$, not only a subsequence; its pair has length at most $n$ and common padding reaches $n$ exactly.  This proves \eqref{eq:min-tv}.  The sample-complexity conclusion follows from Theorem~\ref{thm:main}, whose proof below does not use Theorem~\ref{thm:tv}.  Indeed, $c\log(q^3\log n)\to\infty$ for any fixed admissible $c>0$.
\end{proof}

\subsection{Two-point testing with a strict exponent margin}

For two equally likely hypotheses with observation laws $P,Q$, every randomized decision rule has average success at most $(1+\TV(P,Q))/2$.  If $\varphi$ is the conditional probability of choosing the first hypothesis, its average success is
\[
 \frac12\bigl(1+\E_P\varphi-\E_Q\varphi\bigr)
 \le\frac12(1+\TV(P,Q)),\qquad 0\le\varphi\le1.
\]
This is the standard two-point comparison; trace-reconstruction versions appear in \citet[Appendix A.2]{HoldenLyons2020}.

Fix $0<c<c_+<1/(4\log2)$ and choose $B,\zeta$ so that
\[
 \alpha:=\frac{(1-\zeta)((B-1)/4-.05)}{B\log2}>c_+.
\]
The $j_0$ preliminary blocks have fixed number depending only on $B$, and therefore
\begin{equation}\label{eq:final-exponent}
 d_D=\alpha\log H_{\rm del}+O_B(1).
\end{equation}
For an absolute constant $a>0$ sufficiently small, \eqref{eq:restored-product} with $m\le a n^{d_D}$ is strictly less than $1/3$ eventually: its first term is at most $\sqrt a$, and its second is at most
\[
 2a(D+1)n^{d_D+1-100(d_D+B+1)}=o(1).
\]
Thus the maximal average success for the pair is at most $2/3$ (in fact strictly less for all sufficiently large instances).  Recall from the channel reduction in Section~\ref{sec:introduction} that $\log H_{\rm del}=\log(q^3\log n)+O(1)$ uniformly in $q$.  The strict inequality $\alpha>c_+$ and \eqref{eq:final-exponent} therefore imply
\begin{equation}\label{eq:intermediate-threshold}
 a n^{d_D}\ge n^{c_+\log(q^3\log n)}
\end{equation}
eventually.  Exact worst-case reconstruction with success $2/3$ would distinguish this pair with that success, so this proves the usual $2/3$ lower bound, with the additional margin needed next.

\subsection{Every fixed positive reconstruction success probability}

\begin{proof}[Proof of Theorem~\ref{thm:main}]
Fix $s\in(0,1]$ and $c$ as in the theorem, and retain $c_+>c$ above.  Choose a fixed positive integer $g$ such that $(2/3)^g<s$.  Put $N=\lfloor n/g\rfloor$ and construct the hard pair at length $N$ using the same requested deletion rate $q$ as the original instance.  Its working rate is still $\delta=\min(q,1/2)$.  Since
\[
 q^3\log N\sim q^3\log n\to\infty,
 \qquad \log N=\log n+O_g(1),
\]
the preceding construction and testing bounds apply to each gadget.

The exponent margin is preserved after dividing the length by $g$.  More explicitly,
\[
 \frac{(\log N)\log(q^3\log N)}
      {(\log n)\log(q^3\log n)}\longrightarrow1,
\]
so $c_+>c$ gives
\begin{equation}\label{eq:gadget-margin}
 N^{c_+\log(q^3\log N)}
 \ge n^{c\log(q^3\log n)}
\end{equation}
eventually.  This is why the strict intermediate exponent was chosen before fixing the number of gadgets.

Take $g$ independent fair choices between the two length-$N$ words, concatenate the chosen words, and append a common known tail of length $n-gN$.  This gives $2^g$ distinct length-$n$ hypotheses.  Give the estimator the stronger observation in which every trace is separated at the gadget boundaries, together with an independent trace of the tail.  Concatenation simulates the usual observation exactly.

The amplification of reconstruction difficulty using independent hard blocks also appears in \citet[Section 3.3]{MPV2014} and \citet[Section 4]{HoldenLyons2020}.  For the fixed product prior above, the following posterior calculation justifies the multiplication of optimal success probabilities.

Under the product prior, the separated data from different gadgets are independent.  Their posterior distributions factor as well: for a hypothesis vector $\theta=(\theta_1,\ldots,\theta_g)$ and separated data $z=(z_1,\ldots,z_g)$,
\[
 \Prob(\theta\mid z)=\prod_{j=1}^g\Prob(\theta_j\mid z_j).
\]
Hence the largest conditional probability of a correct joint guess is the product of the largest single-gadget posterior probabilities.  Averaging, independence of the $z_j$ shows that the optimal joint success is the product of the optimal single-gadget average successes.  The tail has the same law under every hypothesis and changes none of these posteriors.

For any positive integer sample budget $m\le n^{c\log(q^3\log n)}$, Equations~\eqref{eq:intermediate-threshold} and~\eqref{eq:gadget-margin} bound every single-gadget success by $2/3$.  The joint average success is therefore at most $(2/3)^g<s$.  When $m=0$, the posterior remains the uniform prior, so the optimal joint success is $2^{-g}\le(2/3)^g<s$.  A reconstruction procedure with worst-case success at least $s$ would have average success at least $s$ under this finite prior, a contradiction even in the stronger experiment.  No nonnegative integer budget at most the claimed threshold suffices, which implies \eqref{eq:main-bound}.
\end{proof}
